\documentclass[11pt]{article}
\usepackage{amsmath,amssymb,amsthm}
\setlength{\emergencystretch}{2em}
\newtheorem{theorem}{Theorem}
\newtheorem{lemma}[theorem]{Lemma}
\newtheorem{corollary}[theorem]{Corollary}
\newcommand{\im}{\operatorname{im}}
\newcommand{\Fix}{\operatorname{Fix}}
\newcommand{\Sp}{\operatorname{Sp}}
\begin{document}
\title{Symplectic flags in string generating sets}
\author{Human Compression}
\date{}
\maketitle

We prove a criterion excluding every string marking of rank at least five
with specified first and last generators. In dimension six the criterion
uses only the ranks and product of their differences from the identity.
It does not require either intermediate group centralizer, the complete
linear commutant of a prefix, or a separate test for each second generator.

The ingredients are standard identities for adjoints, images and
orthogonal complements, together with a dimension count for an involution
in characteristic two. The contribution of this note is their combination
in the endpoint criterion and its proposed application to a retained
finite-group model. No assertion of historical priority is made.
The original named-group identification, representation and endpoint
transport remain imported premises. No computation, operational
acceptance, census, open-question solution or formalization is claimed.

Throughout, \(k\) is a field of characteristic two, \(V\) is a
six-dimensional \(k\)-space, and \(\beta\) is a nondegenerate alternating
bilinear form on \(V\). Matrices act on \emph{column vectors}.
An ordered sequence \(t_0,\ldots,t_{r-1}\) is a string marking if
\(t_i^2=1\) and \([t_i,t_j]=1\) whenever \(|i-j|\geq2\).
The intersection property is not needed below.

\begin{lemma}\label{lem:adjoint}
If \(s\in\Sp(V,\beta)\) satisfies \(s^2=1\), then
\[
 s-1=(s-1)^*,\qquad
 (\im(s-1))^\perp=\ker(s-1).
\]
If, in addition, \(\operatorname{rank}(s-1)=3\), then
\[
       \im(s-1)=\ker(s-1)
\]
is a Lagrangian subspace.
\end{lemma}
\begin{proof}
An isometry has adjoint \(s^*=s^{-1}\). Hence
\((s-1)^*=s^{-1}-1=s-1\).
For any linear map \(T\), nondegeneracy gives
\((\im T)^\perp=\ker T^*\), which proves the second equality.
In characteristic two, \((s-1)^2=s^2-2s+1=0\).
Its image is therefore contained in its kernel. In the rank-three case
both have dimension three, so they are equal. The second equality then
says precisely that this subspace equals its orthogonal complement.
\end{proof}

\begin{lemma}\label{lem:three}
Let \(U\) have dimension three, and let \(B\in\operatorname{GL}(U)\)
satisfy \(B^2=1\). If \(B\ne1\), then
\[
       \operatorname{rank}(B-1)=1,\qquad \dim\Fix_U(B)=2.
\]
\end{lemma}
\begin{proof}
Put \(T=B-1\). Characteristic two gives \(T^2=0\), so
\(\im T\subseteq\ker T\). Rank-nullity implies
\(2\operatorname{rank}T\leq3\). Since \(T\ne0\), its rank is one;
its kernel consequently has dimension two.
\end{proof}

\begin{theorem}\label{thm:endpoints}
Let \(A,E\in\Sp(V,\beta)\) be involutions. Set
\[
       N=A-1,\qquad M=E-1,\qquad
       P=\im N,\qquad U=\ker M.
\]
Suppose
\[
        \operatorname{rank}N=2,\qquad
        \operatorname{rank}M=3,\qquad NM=0.             \tag{1}
\]
Let \(r\geq5\), and let \(t_0,\ldots,t_{r-1}\) be a string marking
in \(\Sp(V,\beta)\), with \(t_0=A\) and \(t_{r-1}=E\).
Then the subgroup generated by the marking stabilizes either \(P\)
or \(U\). Both are proper nonzero subspaces.
\end{theorem}

\begin{proof}
The string relation gives \(AE=EA\), whence \(NM=MN\).
By Lemma~\ref{lem:adjoint}, \(U=\ker M=\im M\) is Lagrangian.
The equation \(MN=0\) gives \(P\subseteq U\), and \(NM=0\) gives
\(U\subseteq\ker N\). In particular, \(A\) fixes \(U\) pointwise.
The dimensions of \(P\) and \(U\) are two and three respectively.

Write \(B=t_1\). Since \(r\geq5\), the string relation gives
\(BE=EB\). Thus \(B\) preserves \(U=\ker(E-1)\).

First suppose that \(B(P)=P\). Every other generator either equals
\(A\) or commutes with \(A\), so it preserves \(\im(A-1)=P\).
The whole generated subgroup therefore stabilizes \(P\).

Now suppose that \(B(P)\ne P\). Choose \(p\in P\) with \(Bp\notin P\).
Then
\[
        v=(B-1)p\in U\setminus P.
\]
In particular, the involution induced by \(B\) on \(U\) is nontrivial.
Lemma~\ref{lem:three} shows that it has a two-dimensional fixed space
and that its difference from the identity has image \(kv\).
Put
\[
       Q=\im(B-1),\qquad
       F=\ker(A-1)\cap\ker(B-1).
\]
Since \(\dim P=2\) and \(v\in U\setminus P\), we have
\[
                   U=P+kv\subseteq P+Q.
\]
By Lemma~\ref{lem:adjoint} applied to \(A\) and \(B\),
\[
 F=P^\perp\cap Q^\perp=(P+Q)^\perp
       \subseteq U^\perp=U.                         \tag{2}
\]
As \(A\) fixes \(U\) pointwise, (2) implies
\[
                       F=\Fix_U(B).
\]
Thus \(\dim F=2\). If \(P\subseteq F\), then \(B\) fixes \(P\)
pointwise, contradicting \(B(P)\ne P\). The two distinct planes
\(P,F\) in \(U\) therefore satisfy
\[
                            P+F=U.                 \tag{3}
\]

Any linear map \(D\) commuting with \(A\) preserves \(P=\im(A-1)\).
If it also commutes with \(B\), it preserves the joint fixed space
\(F\). Equation (3) consequently shows that every such \(D\)
preserves \(U\).

The penultimate generator \(t_{r-2}\) commutes with \(t_0=A\) and
\(t_1=B\): its distance from \(t_1\) is \(r-3\geq2\).
It therefore preserves \(U\). Every other generator either equals
\(E\), or has distance at least two from \(E\); it thus also
preserves \(U=\ker(E-1)\).
This proves that the whole generated subgroup stabilizes \(U\) in
the second case.
\end{proof}

\begin{corollary}\label{cor:proper}
Let \(\rho:G\to\Sp(V,\beta)\) be a representation and let
\(a,e\in G\) have images \(A,E\) satisfying (1).
Suppose that some elements of \(\rho(G)\) fail to preserve \(P\)
and \(U\), respectively. Then no string marking of \(G\) of rank
at least five with first generator \(a\) and last generator \(e\)
generates \(G\).
\end{corollary}
\begin{proof}
The image of every such marking stabilizes one of \(P,U\) by
Theorem~\ref{thm:endpoints}. It therefore omits a member of
\(\rho(G)\). If the original marking generated \(G\), its image
would generate \(\rho(G)\), a contradiction.
Faithfulness is not required.
\end{proof}

A single element failing to preserve both spaces is sufficient.
Alternatively, irreducibility of \(\rho(G)\) implies the escape
hypotheses. The explicit escape version is preferable in a finite
certificate, since it avoids a separate irreducibility claim.

\paragraph{A broader form.}
The same argument works in dimension \(2m\), with \(U\) Lagrangian,
\(\dim P=m-1\), \(U\subseteq\ker(A-1)\), \(P\subseteq U\),
provided that the involution induced by \(B\) on \(U\) has rank at most
one after subtracting the identity. In the second case that rank
is one, so \(\dim F=m-1\) and \(P+F=U\).
Dimension six is useful because the required rank bound on a
three-dimensional involution is automatic.

\section*{The retained six-dimensional matrices}

Let \(e_1,\ldots,e_6\) denote the column basis and let
\(E_{ij}\) be the matrix unit sending \(e_j\) to \(e_i\).
Use the form
\[
             \beta(x,y)=\sum_{i=1}^6 x_i y_{7-i}.       \tag{4}
\]
Its Gram matrix \(J\) has ones on the antidiagonal and zeros elsewhere.

The retained endpoint matrices have the form
\[
 A=I+\mu(E_{15}+E_{26}),\qquad
 E=I+\begin{pmatrix}0&T\\0&0\end{pmatrix},\qquad
 T=\begin{pmatrix}
        1&\alpha&\delta\\
        \gamma&0&\alpha\\
        0&\gamma&1
    \end{pmatrix},                                  \tag{5}
\]
where \(\mu,\delta,\gamma\ne0\).
The block notation in (5) uses the decomposition
\(\langle e_1,e_2,e_3\rangle\oplus
 \langle e_4,e_5,e_6\rangle\).

The matrix \(N=A-I\) has rank two, with
\[
          P=\langle e_1,e_2\rangle.
\]
The matrix \(M=E-I\) is zero on the first summand and maps the second
to the first through \(T\). In characteristic two,
\[
                    \det T=\delta\gamma^2\ne0.
\]
Hence \(M\) has rank three and
\[
       U=\ker M=\im M=\langle e_1,e_2,e_3\rangle.
\]
Both \(N\) and \(M\) map the second summand into the first and
annihilate the first. Thus \(NM=MN=0\).
These calculations establish every matrix condition (1)
without inspecting a second generator.

For completeness, both matrices in (5) preserve (4).
If \(J_3\) is the antidiagonal \(3\) by \(3\) matrix, an upper
unipotent block matrix preserves \(J\) precisely when its upper
right block becomes symmetric after left multiplication by \(J_3\).
For the two blocks in (5), those matrices are
\[
 \begin{pmatrix}0&0&0\\0&0&\mu\\0&\mu&0\end{pmatrix},
 \qquad
 \begin{pmatrix}0&\gamma&1\\\gamma&0&\alpha\\1&\alpha&\delta\end{pmatrix}.
\]
This verifies form preservation of the endpoints. It does not
by itself prove that they are images of the required original
group elements; that correspondence is retained separately.

\paragraph{The acquired field codes.}
In the retained GF(8) exponent codec, zero has code \(0\), one has
code \(1\), and code \(i\), \(1\leq i\leq7\), denotes
\(\zeta^{i-1}\) for the acquired element \(\zeta\) of order seven.
The values in (5) are
\[
    \mu=\text{code }4,\quad \alpha=\text{code }5,\quad
    \delta=\text{code }3,\quad \gamma=\text{code }2.
\]
The nonzero conditions above therefore hold. For example,
code \(2\) does not denote the integer scalar two.

\paragraph{The original ambient generators.}
The two retained ambient matrices are
\[
 D=\operatorname{diag}(\text{code }2,\text{code }7,
       \text{code }3,\text{code }6,\text{code }2,\text{code }7)
\]
and
\[
 C=\begin{pmatrix}
 1&0&1&0&0&0\\
 1&0&0&0&0&0\\
 0&1&0&0&1&0\\
 0&1&0&0&0&0\\
 0&0&0&1&0&1\\
 0&0&0&1&0&0
 \end{pmatrix}.                                      \tag{6}
\]
The three products of paired diagonal entries of \(D\) are one,
so \(D^{\mathsf T}JD=J\).
The columns of \(C\) are
\[
 e_1+e_2,\quad e_3+e_4,\quad e_1,\quad
 e_5+e_6,\quad e_3,\quad e_5.
\]
Their nonzero pairings under (4) are exactly the antidiagonal
pairs, all with value one; the apparent extra pairing between
the first and fourth columns is \(1+1=0\).
Thus \(C^{\mathsf T}JC=J\).

Under the retained premise that \(D,C\) generate the image of the
original group, these two equalities prove that its entire image
is symplectic. No named-group recognition is needed for this step.
Moreover,
\[
                         Ce_2=e_3+e_4
\]
belongs to neither \(P\) nor \(U\).
The same acquired ambient element therefore supplies both escapes
required by Corollary~\ref{cor:proper}.

\paragraph{Consequence for the fixed original pair.}
Subject to the retained original representation, endpoint membership
and ambient-generator correspondences, \emph{every} rank-five string
completion with this endpoint pair fails to generate the original
group. The same conclusion holds for every higher rank.
There is no condition on the second generator beyond the string
relations and group membership. In particular, no second-generator
matrix, full centralizer, commutant basis or search certificate is
needed for this mathematical implication.

The project catalogue associates its family 4 with this fixed
original pair. Its current pending entries with branch labels
\(27,\ldots,36\) have rank five. Applying the mathematical result
to those exact jobs still requires an independently checked,
typed binding from the original catalogue and required-job universe
to this endpoint pair, and an actual acceptance under a suitable
proof interface. This note supplies none of those operational
steps and removes no job from a ledger.

This does not show that all endpoint pairs satisfy (1), and it does not
answer the source paper's Question 4.1. In the retained model, a
generating rank-five example must lie outside the endpoint orbit
treated here. The published constructions are therefore not being
excluded by a change of labels or by an unproved extrapolation.

\paragraph{Automorphisms and action conventions.}
An automorphism of the original group may be applied to an entire
marking before applying \(\rho\). Generation and all string
commutations are preserved. If that automorphism sends the two
endpoints to the retained pair, the theorem applies to the
transformed marking. It need not act linearly in this
representation; no matrix representation of the full automorphism
group is used.

The proof consistently uses columns in the retained matrix group.
If an original action was written on row vectors, this does not
alter the matrix subgroup or its generation question. The proof
does not identify a column with a point of the original permutation
action. In particular, no raw permutation point or contragredient
identification is silently used.

\section*{Finite application obligations}

A small accepting mathematical certificate needs the following
already acquired facts or direct identities.
\begin{enumerate}
\item A representation of the same original group in the retained
      field, and genuine original membership of both endpoints.
\item The exact endpoint matrices (5), with the stated nonzero
      coefficients, in their acquired common basis.
\item Genuine original ambient generators with matrices (6), and
      generation of the image by the supplied ambient list.
\item The acquired field interpretation, the two form identities
      \(D^{\mathsf T}JD=J=C^{\mathsf T}JC\), the rank facts in (1),
      and the two elementary escapes.
\item For a census application, an exact original-job mapping to
      this endpoint pair. Existing classifications and provenance
      keep their original attribution.
\end{enumerate}
All rank and zero-product facts in item 4 have symbolic proofs above.
A verifier may nevertheless check their finite identities afresh.
Any such new verification and its resource accounting are future
work; this note did not run one.

This argument replaces the need to acquire a separate full
\(C_G(a,b)\), a complete \(C_{\operatorname{End}(V)}(A,B)\), or a
proper generated matrix algebra for each second generator.
The earlier commutant theory remains valid and useful for
endpoint pairs failing (1). Neither theory alone completes the
whole F004 census, and no whole-manuscript compression factor is
assigned here.
\end{document}

